

\emph{Pentominous} is a grid logic puzzle based on the twelve pentominoes. A pentomino is a polygon formed by connecting five equal-sized squares edge to edge.

\begin{figure}[!h]
  \centering
  \vspace{-2mm}
  \includegraphics[width=0.6\textwidth]{pentominoes}
  \vspace{-2mm}
  \caption{The twelve pentominoes (up to mirroring and rotations) and their names.}
  \label{fig:pentominoes}
\end{figure}

\vspace{-2mm}
The goal of this puzzle is to divide a grid into regions of size $5$ (that is, pentominoes), so that no two regions that share a side have the same shape.
You are allowed to rotate and mirror the pentominoes, but such rotations and reflections count as the same shape.
The twelve possible shapes can be seen in Figure~\ref{fig:pentominoes}.

In a normal \emph{Pentominous} puzzle, the player is given some pre-filled cells, for which the shape of their region is already predetermined.
In this problem, you are working with a completely blank grid of dimensions $h\times w$, and your task is to create any valid arrangement of pentominoes.

\vspace{-4mm}
